% Scope and conventions are stated in Reading the catalogue.
% Mathematical entries retain their book locators, corrections and proof references.
\documentclass[11pt]{article}
\usepackage[a4paper,margin=22mm]{geometry}
\usepackage{amsmath,amssymb,mathrsfs,longtable,array,bm}
\usepackage[hidelinks]{hyperref}
\usepackage{xurl}
\newcommand{\R}{\mathbf R}
\newcommand{\N}{\mathbf N}
\newcommand{\Leb}{\mathscr L}
\newcommand{\Haus}{\mathscr H}
\newcommand{\Mink}{\mathscr M}
\newcommand{\spt}{\operatorname{spt}}
\newcommand{\dmn}{\operatorname{dmn}}
\newcommand{\Clos}{\operatorname{Clos}}
\newcommand{\Hom}{\operatorname{Hom}}
\newcommand{\card}{\operatorname{card}}
\newcommand{\diam}{\operatorname{diam}}
\newcommand{\im}{\operatorname{im}}
\newcommand{\U}{\mathbf U}
\newcommand{\B}{\mathbf B}
\newcommand{\M}{\mathbf M}
\newcommand{\Lsp}{\mathbf L}
\newcommand{\Lt}{\widetilde{\mathbf L}}
\DeclareUrlCommand\rocq{\def\UrlFont{\footnotesize\ttfamily}}
\newcounter{catalogueentry}
\newcommand{\entry}[3]{%
  \refstepcounter{catalogueentry}%
  \par\addvspace{1.2\baselineskip}%
  \noindent\textbf{#1}\label{entry:\thecatalogueentry}\par\nobreak\smallskip
  \noindent\emph{\entrysecondlabel}\par\nobreak
  \noindent #2\par\smallskip
  \noindent\emph{\entrythirdlabel}\par\nobreak
  \noindent #3\par}
\newenvironment{errors}{%
  \par\begingroup
  \def\entrysecondlabel{Printed passage}%
  \def\entrythirdlabel{Correction and evidence}%
}{\par\endgroup}
\newenvironment{confirmed}{%
  \par\begingroup
  \def\entrysecondlabel{Book page and line}%
  \def\entrythirdlabel{Published correction and status}%
}{\par\endgroup}
\title{\textit{Geometric Measure Theory}\\Errors detected by the formalisation}
\author{Compiled from the Rocq formalisation \texttt{federer-rocq}}
\date{}
\begin{document}
\maketitle
\setcounter{tocdepth}{2}
\tableofcontents
\bigskip

\section*{Reading the catalogue}
\addcontentsline{toc}{section}{Reading the catalogue}

This catalogue lists errors, convention qualifications, and published
corrections to Herbert Federer, \textit{Geometric Measure Theory} (1969).
Each local entry gives the book locator, the printed passage, and the
correction with its evidence. Within each category, entries follow book
order and are grouped by chapter.

Sections~\ref{s:false}--\ref{s:proofs} distinguish false statements,
degenerate cases, convention-dependent qualifications, misprints, and
proof errors absent from Federer's \textit{List of corrections} of May
1987. A convention qualification specifies the interpretation under which
a statement holds; it need not refute the statement under every
interpretation. Section~\ref{s:erratum} records an error in the published
list and an editorial normalization issue. Section~\ref{s:confirmed}
lists published corrections adopted in the formalisation.
Section~\ref{s:further-chapter3} records further published Chapter 3
corrections, including bibliographical changes and corrections outside
the formalised scope. These entries give book pages and lines, followed
by the published correction and its status. The complete published list
is transcribed in \texttt{federer-erratum.tex}.

Book locators refer to the printed Springer reprint of the 1969 edition,
which does not incorporate the 1987 corrections. Items are Federer's
numbered paragraphs; page and line numbers refer to the printed book
unless explicitly assigned to the erratum. Printed passages are close
paraphrases unless quotation marks identify a quotation. Displayed
formulas follow the book. Rocq names identify formal counterexamples,
corrected statements, or supporting results. Arguments outside the
formalisation are marked as such.

Federer's notation is used: $\B(a,r)$ and $\U(a,r)$ are the closed and
open balls, $A\sim B$ is set difference, and $\bigodot_m V$ is the $m$-th
symmetric power. Entries are stated over~$\R$ unless explicitly assigned to the complex
extension of 3.4.12. Where Chapter 1 needs no order, the formalisation
works over a field of characteristic other than two.

The review covers the formalised assertions in Chapters 1 and 2,
Sections 3.1--3.3, items 3.4.1--3.4.12, and items 4.1.1--4.1.27, with the following limits:
\begin{itemize}
\item Chapter 1: item 1.1.3 is deferred. The questions in 1.8.3--1.8.4
are retained as questions.
\item Chapter 2: the book examples in 2.8.8--2.8.9 are excluded;
the Ward construction cited in 2.10.27 is intentionally omitted at the
user's request, and 2.10.46 is partial. The complete comparison and ratio assembly in 2.10.6 is proved. The real-degree extensions listed in
\path{docs/checkpoints/coverage-completion.md} are proved. Source corrections and
convention qualifications remain as recorded in
\path{docs/chapter2/2.10.md}. Item 2.10.49 is a historical source account.
\item Chapter 3: the second paragraph of 3.2.40 (Kneser's example) is
deferred. Item 3.2.45 is complete with its checked quantitative
correction. Items 3.3.11 and 3.3.22 are source accounts. The numbered conclusions of 3.4.8 are validated, including
the compact-component assertion of 3.4.5. The stronger original real
compact-union sentence in the proof of 3.4.8(5) now has a supplementary
proof for arbitrary proper ideals; this is an expanded proof, not a source
correction. Its numbered Borel conclusion has a separate proof. The stronger
complex auxiliary sentence is not asserted by this supplement. The complex
extension in 3.4.12 includes the local covering qualification
and the corrected real/complex dimension comparison below.
\item Chapter 4: items 4.1.1--4.1.27 are included, with the
zero-dimensional nonembedding correction in 4.1.1 below.
Item 4.1.5 retains the original compact scalar-test gauge and its Radon
identification in the finite case; it introduces no new source correction.
Item 4.1.6 reuses the existing degree-zero absolute-scalar convention
for field mass and comass; it introduces no new source correction.
Item 4.1.7 includes the corrected half-line primitive uniqueness
assertion on p.~359, together with actual distinct primitives and the
original support obstruction.
Item 4.1.8 includes the Cartesian product, original measure comparisons,
finite Stokes and the exact ordered scalar-box identity, with the
p.~360 ambient-space misprint corrected below. Its registered build,
complete recursive kernel check and strict dependency audit pass.
Item 4.1.9 includes both current and form homotopy identities, the sharp
variation bound and all three affine mass estimates. It introduces no new
mathematical correction. Its registered build, complete recursive kernel
check and strict dependency audit pass.
Item 4.1.10 includes the original compact-test primitive, all compact-current
cone conclusions and every star-shaped open domain. It applies the published
superscript correction below and introduces no new mathematical correction.
The registered build, complete recursive kernel check and strict dependency
audit pass. Item 4.1.11 includes every original join, mass, simplex, prism, factorial
rectangle and associativity assertion on pp.~364--366. Its algebraic laws
retain arbitrary compact currents, and its mass bounds retain their
original representation hypotheses and exact coefficients. Dimension zero
and degenerate vertices are included. Registered compilation, the complete
recursive kernel check and the strict dependency audit pass. No new
mathematical source error was found. Item 4.1.12 includes the original
flat seminorms, attained mass decomposition, flat-chain closure and
completeness, local criterion, localisation, boundary and scalar laws,
both Cartesian estimates, and both in-book counterexamples on
pp.~367--369. The original partial-operation domains, support indices and
exact coefficients remain explicit. Its registered build, complete
recursive kernel check and strict dependency audit pass. No new
mathematical source error was found. Item 4.1.13 retains the exact
original compact segment hull, both flat bounds for smooth images,
the support-neighbourhood argument and the original variation integrals
and differential weights. Its registered build, complete recursive kernel
check and strict dependency audit pass. No new mathematical source error
was found. Item 4.1.14 retains the original compact and proper local
Lipschitz images, every smooth approximation condition, both homotopy
formulas, continuous Lebesgue-almost-everywhere differential majorants,
support-only affine mass and normal flat bounds, and the class-one action
on every continuous form. Its registered build, complete recursive kernel
check and strict dependency audit pass. No new mathematical source error
was found. Item 4.1.15 retains the original radial-map and smoothing
proof of support coincidence, retraction and point-support formulas,
the three section-image sets, the affine-plane obstruction and the
actual compact finite-mass nonflat example with its bounded-derivative
smooth approximation. The final sentence on p.~373 has the local
degree correction recorded below. Its registered build, complete
recursive kernel check and strict dependency audit pass.
Item 4.1.16 retains the original nearest-point retraction onto every
nonempty closed convex set, its quadratic and variational formulas,
the Lipschitz bound one, the exact compact convex hull of each flat
current's support, and the equality of flat norms at every larger
compact index. Empty support and zero ambient dimension are included.
The published ambient-space correction on p.~373 is recorded below;
no new source error was found. The registered build, complete
recursive kernel check and strict dependency audit pass.
Item 4.1.17 retains the exact normal mass-closure equality and
contraction by every original variation-summable differential form
with the same compact flat index. Representation and local
summability are derived from the original finite-mass hypotheses;
the coefficient need not be bounded, smooth or compactly supported.
The proof adopts the published normal approximant $Q+\partial S$,
erratum p.~7, book p.~374, line~16. No new local source error was
found. Registered compilation, complete recursive kernel checking
and the strict dependency audit pass.
Item 4.1.18 retains the original reflected translation average,
the smoothing deformation current, every current and form homotopy
case, and the sharp mass, support and compact flat-error bounds,
including infinite masses. The actual averaged translated Radon
variation bounds every set and gives weak variation convergence
without finite total mass or compact support of the represented
current. Both original compact $L^1$ pair propositions retain the
exact enlarged indices and cost bounds on every open $U$. The
normal approximation, residual iteration, common half-radius support,
both infinite sums and downward degree induction are proved without
assuming finite mass or normality of the initial flat current.
The opaque onto map, quotient fibres, top-degree characterization
and both inclusions in the literal $J$-kernel closure are included.
The raw integral cost is a seminorm under null-set changes. The
$L^1$ convergence laws respect almost equality; the literal kernel
closure keeps the original raw pair space and $J$ topology.
This carrier distinction is not a new error claim. Empty supports
and all original degree cases remain. No new source error was
found. Registered compilation, the complete recursive kernel check
and the strict dependency audit pass.
Item 4.1.19 includes the original common-bound flat cochains,
their least nonnegative norm and actual coboundary, the compact
$L^1$ pair functional and signed unit-ball supremum, both measurable
form actions, almost uniqueness, the weak equation, the unique converse,
the class-one and interior smoothing criteria, and the coboundary pair
$(\psi,0)$. The literal global-comass converse and pointwise class-one
parenthesis have the two local representative corrections recorded below.
Selected representatives retain the original global formula on nonempty
$U$; the empty-domain version adjoins the nonnegative threshold zero.
The original current union, compact indices, arbitrary open domains and
all natural degrees remain. The final Whitney comparison and historical
continuation are source commentary. Registered compilation, the complete
recursive kernel check and the strict dependency audit pass.
Item 4.1.20 includes the original positive-degree flat-current theorem:
$\mathcal I_1^m(\operatorname{spt}T)=0$ implies $T=0$ on every original
open $U$. The proof retains the original orthogonal projection, finite
coordinate and top-degree density arguments. No normality, finite-mass or
representation premise is added. No new source correction is introduced.
Registered compilation, the complete recursive kernel check and the strict
dependency audit pass.
Item 4.1.21 includes both original positive-degree null-set corollaries.
The first retains arbitrary $B$ and the relative closedness of
$\operatorname{spt}T\setminus B$, without a finite-mass or representation
premise. The second gives zero variation to every arbitrary
$\mathcal I_1^m$-null subset of the original $U$ when $T$ is flat and
$M(T)<\infty$. A Borel null hull precedes Radon inner regularity;
compact measurable restrictions retain exact support, flatness and mass.
No new source correction is introduced. Registered compilation, the
complete recursive kernel check and the strict dependency audit pass.
Item 4.1.22 includes the fixed-index integral polyhedral subgroups,
their real algebraic spans, and the two unions over every original
compact subset of the open domain. Independent generation laws and
opaque carriers use the existing oriented simplex currents with full
vertex hulls in the prescribed index. The algebra, fixed-index support,
normality, flatness, boundary and open restriction laws are included.
Empty indices and signed degree-zero points retain their original
meaning. No new source correction is introduced. Registered compilation,
the complete recursive kernel check and the strict dependency audit pass.
Item 4.1.23 includes real polyhedral approximation with the original
compact flat error and mass error, all five printed reductions, and the
full transverse atomic construction on pp. 379--380. The original
compact target index need not be convex, and the input mass may be
infinite. Every generator hull stays in the prescribed index. Signed
degree-zero terms and the transverse one-point space in top degree are
retained. The normal-mass strengthening remains a later obligation of
4.2.24. No new source correction is introduced. Registered compilation,
the complete recursive kernel check and the strict dependency audit pass.
Item 4.1.24 includes the rectifiable currents defined by mass
approximation with Lipschitz images of integral polyhedral chains, the
integral currents, integral flat chains, their compact unions and local
classes, every inclusion of the diagram on p.~381, the integral flat
size and its complete metric, mass completeness and density, the series
estimate, and the ordinary flat finite-mass decomposition with an attained
flat filling. The equality $\mathscr R_0=\mathscr P_0$ retains integer
coefficients. No new source correction is introduced. Registered
compilation, the complete recursive kernel check and the strict
dependency audit pass.
Item 4.1.25 includes the Lebesgue action of the Lipschitz image of a
compactly supported summable current, the full-fibre density $\eta$ with
its measurability and summability, the identity
$f_\#(\Leb^m\wedge\xi)=\Haus^m\wedge\eta$, the associated tangent and
the signed full-fibre formula. The simplicity of $\eta(y)$ is qualified
in degree zero below. No other source correction is introduced.
Registered compilation, the complete recursive kernel check and the
strict dependency audit pass.
Item 4.1.26 includes the signed degree of a Lipschitzian map
$f:\mathbf R^m\to\mathbf R^m$ on a bounded $\Leb^m$ measurable set $A$,
the identity $f_\#(\mathbf E^m\mathop{\llcorner}A)=\mathbf E^m\mathop{\llcorner}(\operatorname{degree}f|A)$,
its almost constancy on each component of
$\mathbf R^m\sim f(\operatorname{Bdry}A)$, and both consequences for
connected $A$ and univalent $f|A$. It introduces no source correction.
Registered compilation, the complete recursive kernel check and the
strict dependency audit pass.
Item 4.1.27 includes the prefix-product map $f$ of the open unit cube $C$,
the exact image $f(C)$, the wedge formula for the partial derivatives at
every point of $\R^m$, the identity
$[u_0,\dots,u_m]=f_\#(\mathbf E^m\mathop{\llcorner}C)$, and
$[u_0,\dots,u_m]=\Haus^m\wedge\eta$. It uses the published correction of
p.~384, line~19, below. The univalence clause is qualified in the
degenerate case $n+1<m$ below. Registered compilation, the complete
recursive kernel check and the strict dependency audit pass.
Items 4.1.28--4.1.34 remain pending.
Chapter 5 is outside the review.
\end{itemize}
An entry that records a published correction outside this scope does not
claim a formal proof of the corrected assertion. Absence from the
catalogue does not certify a passage as error-free.

The accompanying \texttt{docs/divergences.md} records changes of statement,
convention, and proof; the section guides under \texttt{docs/chapter1},
\texttt{docs/chapter2}, \texttt{docs/chapter3}, and \texttt{docs/chapter4} give detailed coverage
and Rocq declarations. The documentation index is \texttt{docs/toc.md},
and source editions are identified in \texttt{docs/bibliography.md}.
The project's GitHub mirror is at
\url{https://github.com/sk197876/math/tree/main/federer1969}.

\section{False statements}\label{s:false}
\subsection{Chapter 1}

\begin{errors}
\entry{1.4.4, p.~20}{For every associative algebra $W$ with unit, the
identity of $\Hom(V,W)$ extends to a homomorphism of the exterior algebra
of $\Hom(V,W)$ over $\R$ into the algebra of alternating $W$-valued forms.}
{False for noncommutative $W$: such an extension forces all values of a
one-form to commute. Take $V=\R^2$, $W$ the free associative algebra on two
letters $X,Y$, and $\alpha(e_1)=X$, $\alpha(e_2)=Y$; then
$(\alpha\wedge\alpha)(e_1,e_2)=XY-YX\neq0$. \emph{Correction:} require $W$
to be commutative; the homomorphism then exists and is unique.\newline
\rocq{noncommutative_exterior_extension_impossible}, \rocq{Omega_W},
\rocq{Omega_W_unique}}
\entry{1.4.4, p.~20}{The homomorphism of the preceding entry is surjective
whenever $\dim V<\infty$.}
{False for every $V$ unless $W=\R\cdot1_W$: in degree zero the map is
$c\mapsto c\cdot1_W$. For $W=\R\times\R$ the constant zero-form $(1,0)$ has
no preimage. \emph{Correction:} delete the assertion. In finite dimension
the isomorphism holds for the exterior algebra over a commutative $W$ of
the $W$-module $\Hom(V,W)$, as at the end of the same parenthesis.\newline
\rocq{coefficient_pair_not_surjective},
\rocq{Omega_W_surjective_requires_scalar_coefficients},
\rocq{Omega_relative_finite_bijective}}
\end{errors}

\subsection{Chapter 2}

\begin{errors}
\entry{2.4.12, p.~86}{Corollary of the Hahn--Banach theorem: for a
seminorm $\sigma$ on $V$ with values in $\{t:0\le t\le\infty\}$ and
$v\in V$ there is a linear $\alpha:V\to\R$ with $\alpha\le\sigma$ and
$\alpha(v)=\sigma(v)$.}
{False when $\sigma(v)=\infty$, since $\alpha$ is real valued. On $V=\R$
take $\sigma(0)=0$ and $\sigma(x)=\infty$ for $x\neq0$, and $v=1$.
\emph{Correction:} assume $\sigma(v)<\infty$.\newline
\rocq{infinite_seminorm_has_no_norming_functional},
\rocq{extended_seminorm_norming_functional}}
\entry{2.5.18(2), p.~112}{For $f,g,h:\R\to\R$ with $f,g$ continuous and
$g$ increasing, $\int_a^b(f\circ g)\,d(h\circ g)=\int_{g(a)}^{g(b)}f\,dh$;
``therefore $g_\#$ maps the measures induced by $h\circ g$ onto the measures
induced by~$h$''.}
{False unless $\im g=\R$. For $g(x)=x/(1+|x|)$ and $h(x)=x$, the image under
$g$ of the measure induced by $h\circ g$ vanishes on $\{t:2\le t\le3\}$,
whereas the measure induced by $h$ is $\Leb^1$.
The integral identity is unaffected. \emph{Correction:} assume
$\im g=\R$ in the final sentence.\newline
\rocq{stieltjes_image_unrestricted_false},
\rocq{stieltjes_image_requires_onto},
\rocq{stieltjes_image_induced_measures}}
\entry{2.6.5, p.~119}{The product of measures is associative, and
$\alpha_1\times\cdots\times\alpha_n$ does not depend on the association.}
{False for arbitrary measures. Let $\alpha=\beta=\Leb^1$ and let $\gamma$ be
a measure over a nonempty $Z$ with $\gamma(A)=\infty$ for every nonempty
$A\subset Z$. The product of the diagonal of $[0,1]$ and $Z$ has measure $0$
in $(\alpha\times\beta)\times\gamma$ and $\infty$ in
$\alpha\times(\beta\times\gamma)$. \emph{Correction:} assume that each
factor space is the union of countably many measurable sets of finite
measure (for the triple product, the first and third factors suffice).\newline
\rocq{lebesgue_product_not_associative},
\rocq{product_measure_associative_countable},
\rocq{finite_product_association_invariance}}
\entry{2.7.16(4), p.~138}{Weil's condition is violated by every
$t\in GL(n,\R)^W$ with $\det t=1$ and $\det(t|W)\neq1$.}
{The condition reads $\chi(t)\,|\det(t|W)|^{-n}\,|\det t|^m=1$, so it holds
when $\det(t|W)=-1$. \emph{Correction:} replace $\det(t|W)\neq1$ by
$|\det(t|W)|\neq1$. For $0<m<n$ the map acting as $2^{n-m}$ on $W$ and as
$2^{-m}$ on a complement has $\det t=1$ and $|\det(t|W)|=2^{m(n-m)}$, so
the conclusion stands.\newline
\rocq{grassmann_no_covariant_integral},
\rocq{stabiliser_scaling_exists}}
\entry{2.10.1, p.~170}{Alternate description: $\psi(A)$ is the infimum of
all $t$ such that every open covering $H$ of $A$ has a countable subfamily
$G$ of $F$ covering $A$, with each member in some member of $H$ and
$\sum_{S\in G}\zeta(S)<t$. ``For any topological space $X$, it remains
clear that $\psi$ measures $X$.''}
{False: the formula need not be monotone. Let $X=\{p,q\}$ with open sets
$\varnothing$, $\{q\}$, $X$, let $F=\{X\}$ and $\zeta(X)=1$. The value is
$\infty$ on $\{q\}$ and $1$ on $X$. \emph{Correction:} quantify over open
covers of the whole space $X$. The result is a measure, every continuous map
into a metric space is measurable, and it agrees with the printed formula
when $X$ is a metric space.\newline
\rocq{printed_topological_refinement_not_measure},
\rocq{topological_caratheodory_is_measure},
\rocq{topological_caratheodory_continuous_measurable},
\rocq{topological_caratheodory_metric_agreement}}
\entry{2.10.11, p.~176}{The parenthetical sentence says that the
Hausdorff measure in the Lipschitz multiplicity argument may be replaced
by spherical measure.}
{False for arbitrary metric targets. Let $X=\{0,1\}^{\N}$ and set
$d(x,y)=2^{-j}$ when $j$ is the first differing coordinate, with
$d(x,x)=0$. This metric space is compact, complete and separable.
Give $Y=X\times\R$ the maximum metric and put $f(x)=(x,0)$. Then
\[
 \operatorname{Lip}f=1,\qquad
 \mathscr S_X^1(X)=1,\qquad \mathscr S_Y^1(f(X))=2.
\]
The image is compact and Borel; its injective multiplicity has integral
equal to its spherical mass. Thus the original hypotheses hold, but the
printed bound would require $2\le1$.
\emph{Correction:} for arbitrary metric spaces the factor $2^p$ gives
\[
 \int N(f|A,y)\,d\mathscr S^p(y)
 \le2^p(\operatorname{Lip}f)^p\mathscr S^p(A).
\]
This follows from $\mathscr H^p\le\mathscr S^p\le2^p\mathscr H^p$ and
the original Hausdorff image bound, for every real $p\ge0$. The extra
factor two is sharp in degree one. This is a project counterexample and
correction, absent from the published erratum.\newline
\rocq{federer_2_10_11_spherical_counterexample},
\rocq{binary_vertical_spherical_values},
\rocq{federer_2_10_11_real_spherical_corrected}}
\entry{2.10.17(2), p.~179}{If $\mu$ is regular, $A\subset X$, $t>0$ and
$0\le(C)\limsup_{S\to x}\mu(A\cap S)/\zeta(S)<t$ whenever $x\in A$, then
$\mu(A)\le t\,\psi(A)$.}
{False. By 2.1.1 the quotient is undefined when $\zeta(S)=0$, and by
2.8.16 (p.~151) such $S$ are omitted from $(C)\limsup$. On $\R$ let every set belong to $F$,
$\zeta(S)=0$ for $\card S\le1$ and $\zeta(S)=1$ otherwise, $\mu$ the unit
mass at $0$, $A=\{0\}$. The upper density at $0$ is $1<2$, but
$\mu(A)=1>2\,\psi(A)=0$. \emph{Correction for (2):} set $0/0=0$ and
$a/0=\infty$ for $a>0$ in its auxiliary density quotient; the density at
$0$ is then $\infty$, and the corrected (2) is proved. The same convention
also gives a proved auxiliary version of (3).
\emph{Original (3) is valid and formalised:} its numerator is the uncut
covering measure, not an arbitrary $\mu$. The singleton-generator cover
gives $\varphi_\infty(S)\le\zeta(S)$, so a zero-cost generator has zero
restricted uncut mass. Filling just these undefined quotients with zero
preserves the strict bound below one, and the existing null-locus proof
applies. No correction of the original (3), no additional hypothesis,
and no extension of global partial division is needed.\newline
\rocq{printed_caratheodory_density_bound_false},
\rocq{caratheodory_density_divide}, \rocq{federer_2_10_17_2},
\rocq{federer_2_10_17_3}, \rocq{federer_2_10_17_3_original}}
\end{errors}

\subsection{Chapter 3}

\begin{errors}
\entry{3.1.17, p.~229}{``The converse of Theorem 3.1.16 is clearly also
true'': if for each $\varepsilon>0$ some map $g:\R^m\to\R^n$ of class $1$
satisfies $\Leb^m(A\sim\{x:f(x)=g(x)\})<\varepsilon$, then
$\operatorname{ap}\limsup_{x\to a}|f(x)-f(a)|/|x-a|<\infty$ for $\Leb^m$
almost all $a$ in~$A$.}
{False for nonmeasurable $A$: a constant map on a nonmeasurable subset of
$\R^m$ (2.2.4) has a smooth extension. Finite approximate upper growth
almost everywhere forces $A$ to be $\Leb^m$ measurable (via 3.1.8). The
counterexample uses the definition of 2.9.12 as corrected in Federer's list
(book p.~159 line~7), under which points outside $A$ count against the
density; under the original wording of p.~159 the example satisfies the
conclusion. \emph{Correction:} for $m>0$, finite approximate upper growth at almost all
points of $A$ is equivalent to $\Leb^m$ measurability of $A$ together with
the approximation property.\newline
\rocq{nonmeasurable_domain_converse_counterexample},
\rocq{approximate_growth_ae_domain_measurable}, \rocq{federer_3_1_17}}
\entry{3.2.27, p.~262}{If $W$ is an $(\Haus^m,m)$ rectifiable subset of
$\R^n$, $m\le n$ and $a_\lambda=\int N(\mathbf p_\lambda|W,y)\,d\Leb^my$ for
$\lambda\in\Lambda(n,m)$, then
$[\sum_\lambda(a_\lambda)^2]^{1/2}\le\Haus^m(W)\le\sum_\lambda a_\lambda$.}
{The upper bound is false without measurability of $W$ (3.2.14 does not
require it), even when the integrals exist. Take a
nonmeasurable $V\subset[0,1]$, put $V$ on the segment at height $0$ and
$[0,1]\sim V$ on the segment at height $1$. The coordinate integrals are
$1$ and $0$, but $\Haus^1(W)>1$. \emph{Correction:} assume $W$ is
$\Haus^m$ measurable for the upper bound. The lower bound holds as printed
whenever the integrals exist.\newline
\rocq{rectifiable_coordinate_upper_requires_measurability},
\rocq{rectifiable_coordinate_projection_lower}, \rocq{federer_3_2_27}}
\entry{3.2.37, p.~274}{The upper and lower Minkowski contents
$\Mink^{*m}$ and $\Mink_*^m$ are not measures; the preceding definition
permits every integer $0\le m\le n$.}
{False for $m=0$. Both contents are exactly counting measure on every
subset of $\R^n$, including infinite sets and the positive atom in
$\R^0$. \emph{Correction:} restrict the non-measure assertion to
$1\le m\le n$. For these positive degrees, finite sets have zero content
while an actual countable dense set has infinite content, so countable
additivity fails.\newline
\rocq{zero_minkowski_contents_counting},
\rocq{zero_minkowski_contents_are_measures},
\rocq{positive_minkowski_contents_not_measures}}
\entry{3.2.45, p.~279}{For $0\le u<v\le1$, the normalized
measure ratio is bounded above by
\[
 \frac{\beta_{1/u}(n,m)\,\mathscr Q_{1/u}^m(A)}
      {\beta_{1/v}(n,m)\,\mathscr Q_{1/v}^m(A)}
 \le \left(\frac{\gamma}{\beta_1(n,m)}\right)^{v-u}.
\]}
{The upper bound is false. Take $n=3$, $m=1$, and an actual unit
coordinate segment $A$. The rectifiable equality in 3.2.26 gives
$\mathscr Q_t^1(A)=\Haus^1(A)=1$ for every original exponent.
The actual projection moments satisfy
$\beta_t(3,1)^t=1/(t+1)$ for every finite $t\ge1$.
The displayed coarse choice is $\gamma=\alpha(1)2^{-1}1!=1$.
At $u=0$, $v=1/2$, the claimed bound would require
$\sqrt3\le\sqrt2$.
Nor does excluding $u=0$ suffice: at $u=1/16$, $v=1/8$, the
sixteenth power of the actual ratio is $81/17>2$, whereas the
claimed upper bound has sixteenth power two.
In fact no finite $\gamma>0$ repairs the estimate. Choose
$t=2\gamma+1$, $u=0$, $v=1/t$. The $t$th powers of the actual
ratio and claimed upper bound are $t+1$ and $2\gamma$, respectively.
A valid consequence of the stated convexity and monotonicity is
\[
 0\le h(u)-h(v)\le C\frac{v-u}{v},
 \qquad C\ge h(0)-h(1),\quad 0\le u<v\le1.
\]
This supporting bound alone does not establish the infinity endpoint.
The complete original continuity theorem is now proved for every set,
with both endpoints and zero or infinite masses. A separate uniform
relative projection-volume modulus supplies the infinity argument;
the corrected convex secant supplies the finite-exponent argument.
The product upper inequality and sharp $Q_\infty\le Q_1$ comparison
are proved separately. The full original bounded-convex section
inequality has a finite ordinary integral. The all-set upper section
integral and measurable-first-factor product equality are also proved.
The equality retains the ordinary defined-product domain. The particular
Steiner coefficient comparison follows from individual coordinate-stretch
bounds and their actual anisotropic limits, without assuming a general
mixed-volume theorem.
The false quantitative estimate is a project correction, absent from
the published erratum.\newline
\rocq{unit_segment_convex_exponent_values},
\rocq{projection_three_one_moment},
\rocq{convex_exponent_secant_counterexample},
\rocq{convex_exponent_no_finite_secant_constant},
\rocq{convex_nonincreasing_reciprocal_secant},
\rocq{unit_segment_positive_reciprocal_regression},
\rocq{convex_integral_geometric_measure_continuous},
\rocq{gillespie_euclidean_product_upper},
\rocq{convex_integral_geometric_infinity_below_first},
\rocq{bounded_convex_intrinsic_section_comparison},
\rocq{convex_integral_geometric_section_upper_integral},
\rocq{federer_3_2_45_product}}

\entry{3.2.46, p.~281}{On an arbitrary open $A\subset\R^n$ with a
continuous field of inner products, the infimum of the lengths of
Lipschitzian curves joining $x,y\in A$ is a metric inducing the usual
topology of $A$.}
{The distance need not be finite. Take $A=\R\sim\{0\}$ with the ordinary
inner product. Every continuous curve from $-1$ to $1$ meets $0$, so
there is no admissible curve in $A$ and $\rho(-1,1)=\inf\varnothing=\infty$.
\emph{Correction:} allow an extended distance on $A$, or localize the
finite-metric assertion to a connected component. The formalisation
retains the original extended $\rho$; $\min\{1,\rho\}$ is a finite carrier
used for topology, with agreement of the original cutoffs below one.
The measure and transformation statements retain the original distance.\newline
\rocq{riemannian_punctured_line_infinite_distance},
\rocq{riemannian_distance_no_curve}}
\end{errors}

\subsection{Chapter 4}

\begin{errors}
\entry{4.1.7, p.~359}{In the half-line example, the only
$S\in\mathscr D_1(\R)$ with
$\partial S=f_\#(\partial T)=\delta_b^{\R}$ is given by
\[
 S(\psi)=\int_{-\infty}^{b}\langle 1,\psi(x)\rangle\,d\Leb^1 x.
\]}
{The boundary equation alone does not determine $S$. For $0<b<\infty$,
the two actual currents
\[
 S_-(\psi)=\int_{-\infty}^{b}\langle 1,\psi(x)\rangle\,d\Leb^1 x,
 \qquad
 S_+(\psi)=-\int_b^{\infty}\langle 1,\psi(x)\rangle\,d\Leb^1 x
\]
have boundary $\delta_b^{\R}$ and satisfy
$S_--S_+=\mathbf E^1\ne0$.
\emph{Correction:} add the intended agreement $S\llcorner U=T$ on
$U=(0,\infty)$ to the uniqueness assertion. The preceding constancy
theorem then proves that $S=S_-$. The example's support obstruction
remains valid: no whole-line current with this boundary has support
contained in $f(\spt T)=(0,b]$. This correction is absent from the
published erratum.\newline
\rocq{halfline_primitive_nonuniqueness},
\rocq{halfline_primitive_unique_with_agreement},
\rocq{halfline_current_supported_primitive_impossible}}
\entry{4.1.19, p.~377, converse norm formula}{Every bounded measurable
pair $\phi,\psi$ satisfying the displayed weak equation represents a flat
cochain $\alpha$ with
$\mathbf F(\alpha)=\max\{\mathbf M(\phi),\mathbf M(\psi)\}$,
where $\mathbf M$ is the pointwise global comass of 4.1.7.}
{Take $U=\R$, $m=0$, $\phi=\mathbf1_{\{0\}}$ and $\psi=0$.
Both original coefficient actions and the weak equation hold for the
actual zero cochain: the singleton is Lebesgue-null. Its cochain norm is
zero, while $\mathbf M(\phi)=1$ and $\mathbf M(\psi)=0$.
\emph{Correction:} for arbitrary representatives use the maximum of the
two essential comasses. Suitable null modifications retain the printed
global formula for selected representatives on nonempty $U$. On empty
$U$, both literal global comasses are $-\infty$ and the cochain norm is
zero; the all-domain selected formula is
$\mathbf F(\alpha)=\max\{0,\mathbf M(\phi),\mathbf M(\psi)\}$.
The original existence assertion for selected representatives remains
valid on nonempty $U$; the counterexample refutes only the unqualified
converse norm claim. This correction is absent from the published
erratum.\newline
\rocq{flat_zero_cochain_null_spike},
\rocq{flat_cochain_original_form_essential_norm},
\rocq{flat_cochain_original_selected_global_comass},
\rocq{flat_cochain_selected_nonnegative_global_comass}}
\entry{4.1.19, p.~377, class-one parenthesis}{For $\phi$ of class one,
the weak equation holds if and only if $d\phi=\psi$, for an arbitrary
bounded measurable second form $\psi$.}
{A literal pointwise reading fails. On $U=\R$, take $\phi=0$ and
$\psi=\mathbf1_{\{0\}}\,dx$. The first form is smooth, both forms are
bounded measurable and the weak equation holds, but the actual exterior
derivative graph fails at zero. \emph{Correction:} read $d\phi=\psi$
almost everywhere. If $\psi$ is continuous, this gives pointwise equality
on $U$. The formalisation preserves both versions and adds no continuity
premise to the measurable theorem. This qualification is not stated in
the published erratum.\newline
\rocq{class_one_weak_exterior_null_spike},
\rocq{current_class_one_weak_exterior_graph_iff},
\rocq{current_class_one_weak_exterior_pointwise_iff}}
\end{errors}

\section{Degenerate cases}\label{s:degenerate}
\noindent The literal definitions make the empty product $1$ the only simple
$0$-vector. Several assertions of chapter 1 therefore fail in degree zero.
Other assertions fail for zero spaces or empty sets.
\subsection{Chapter 1}

\begin{errors}
\entry{1.4.4, p.~20}{The homomorphism of the first entry of
section~\ref{s:false} is not injective unless $\dim W=1$.}
{False for $V=0$ and $W=\R\times\R$: the map is $c\mapsto(c,c)$, which is
injective, although $\dim W=2$. \emph{Correction:} assume $V\neq0$ (and $W$
commutative). The corrected assertion is argued in \texttt{docs} by an
independent pair of one-forms and is formalised in full generality:
injectivity forces the coefficient unit to be a singleton basis. Zero and
infinite coefficient dimensions are included. The zero coefficient algebra
is handled by its vanishing unit.\newline
\rocq{zero_domain_Omega_pair_injective},
\rocq{Omega_W_not_injective},
\rocq{Omega_W_injective_requires_coefficient_line},
\rocq{Omega_W_noninjective_unless_coefficient_line}}
\entry{1.4.5, p.~20}{The ordinary tensor product of the exterior algebras of
$V$ and $W$ is neither commutative nor anticommutative (unless $V\cong\R$
or $W\cong\R$).}
{False for zero spaces. The sharp exceptions are $V=0$, $W=0$, or
$V\cong\R\cong W$. \emph{Complete corrected classification (formalised):}
the ordinary product is commutative exactly when both original spaces
have dimension at most one; it is graded anticommutative exactly when
at least one original space is zero. Hence it has neither property
exactly when both spaces are nonzero and at least one contains two
independent vectors. These equivalences impose no finite-dimensional
hypothesis. The first factor remains the full scalar alternating algebra.
The proof reuses the earlier grading, exterior dimension, independent-dual
and commutator laws. Concrete zero-plane, plane-zero, line-plane and
two-line regressions are checked; the earlier $V=W=0$ counterexample
remains.\newline
\rocq{ordinary_tensor_commutative_iff},
\rocq{ordinary_tensor_anticommutative_iff},
\rocq{ordinary_tensor_neither_iff},
\rocq{zero_left_plane_ordinary_exception},
\rocq{line_plane_ordinary_neither}}
\entry{1.6.1, p.~24}{A nonzero $m$-vector $\xi$ is simple if and only if
its associated subspace has dimension~$m$.}
{The reverse implication fails for $m=0$: the scalar $2$ has associated
subspace $\{0\}$ but is not simple. \emph{Correction:} assume $m>0$; for
$m=0$ only the forward implication holds.\newline
\rocq{simple_dimension_zero_counterexample}, \rocq{simple_iff_dimension}}
\entry{1.6.3, p.~24}{For $0\neq\phi\in\bigwedge^mV$ the associated
subspace $\bigwedge^1V\cap\{\alpha:\phi\wedge\alpha=0\}$ has dimension
$\le m$, with equality if and only if $\phi$ is simple.}
{Fails for $m=0$, with the covector $2$. \emph{Correction:} assume
$m>0$.\newline
\rocq{covector_dimension_zero_counterexample},
\rocq{simple_covector_iff_dimension}}
\entry{1.6.3, p.~24}{The isomorphisms $\mathbf D_p$ and $\mathbf D^p$ of
1.5.2 preserve simplicity.}
{Fails when the output degree is $0$, that is $p=n$: for $\dim V=1$,
$\mathbf D_1$ maps twice a basis vector, and $\mathbf D^1$ maps twice the
dual covector, to the scalar $2$. \emph{Correction:} assume $p<n$.\newline
\rocq{D_p_zero_degree_counterexample},
\rocq{D_upper_zero_degree_counterexample}, \rocq{D_p_preserves_simple}}
\entry{1.7.5, p.~32}{Equality holds in $(v_1\wedge\cdots\wedge v_m)\bullet
(v_1\wedge\cdots\wedge v_m)\le\prod_{i=1}^mv_i\bullet v_i$ if and only if $v_1,\ldots,v_m$ are mutually
orthogonal.}
{Fails for $(0,u,u)$ with $|u|=1$. \emph{Correction:} equality holds if and
only if some $v_i=0$ or $v_1,\ldots,v_m$ are mutually orthogonal.\newline
\rocq{hadamard_zero_factor_counterexample},
\rocq{exterior_product_norm_equality}}
\entry{1.8.1, p.~38}{The comass $\|\phi\|=\sup\{\langle\xi,\phi\rangle:
\xi\in\bigwedge_mV,\ \xi$ is simple, $|\xi|\le1\}$ satisfies
$|\phi|\ge\|\phi\|\ge\binom{\dim V}m^{-1/2}|\phi|$, and comass and mass are
dual norms.}
{Fails for $m=0$: the only simple $0$-vector is $1$, so $\|c\|=c$, and
$\|-1\|=-1$. \emph{Correction:} state comass,
mass and the norm assertions for $m>0$.\newline
\rocq{literal_zeroth_comass}, \rocq{literal_zeroth_comass_negative}}
\entry{1.8.1, p.~38}{The mass unit ball $C$ consists of all sums
$\sum_{i=1}^Nc_i\xi_i$ with $\xi_i\in S=\bigwedge_mV\cap\{\xi:\xi$ is simple
and $|\xi|\le1\}$, $c_i>0$, $\sum_{i=1}^Nc_i=1$ and
$N\le\dim\bigwedge_mV=\binom{\dim V}m$.}
{Fails when $\bigwedge_mV=0$, that is $m>\dim V$: then $N=0$ and the
coefficients cannot sum to $1$. \emph{Correction:} assume $m\le\dim V$. The
convex hull equality and the optimal decomposition hold for all $m$.\newline
\rocq{zero_dimensional_positive_convex_impossible},
\rocq{mass_positive_convex_representation}}
\entry{1.10.5, p.~48}{The mass $\|\xi\|$ of $\xi\in\bigodot_mV$ equals the
infimum of the sums $\sum_{j=1}^N|v_{1,j}|\cdots|v_{m,j}|$ over
representations $\xi=\sum_{j=1}^Nv_{1,j}\odot\cdots\odot v_{m,j}$; if
$\dim V<\infty$ one can take $N\le\dim\bigodot_mV$.}
{Fails for $m=0$: every empty product is $1$, so no finite sum of empty
products is $-1$, and the infimum for $\xi=-1$ is $\inf\varnothing=\infty$. \emph{Correction:} assume $m>0$; in degree zero the mass
of $c$ is $|c|$.\newline
\rocq{symmetric_negative_unit_not_empty_sum},
\rocq{symmetric_degree_zero_mass_specification}}
\end{errors}

\subsection{Chapter 2}

\begin{errors}
\entry{2.2.15, p.~71}{For an arbitrary separable metric space $Y$, every
Borel function mapping any metric space into $Y$ is a Baire function if and
only if for every finite subset $Z$ of $Y$ and every $\varepsilon>0$, $Z$ is
contained in the $\varepsilon$ neighbourhood of the image of some
continuous map of $\R$ into~$Y$.}
{Fails for $Y=\varnothing$: every map into $\varnothing$ is a Baire map,
whereas the condition for $Z=\varnothing$ demands a curve $\R\to
\varnothing$. \emph{Correction:} replace the condition by \emph{$Y$ is
empty or} the printed condition.\newline
\rocq{empty_target_borel_baire_counterexample},
\rocq{borel_baire_target_criterion_with_empty}}
\entry{2.3.2, p.~73}{A $\phi$ measurable function $f$, whose domain
contains $\phi$ almost all of $X$, can always be extended to a $\phi$ almost
equal function with domain~$X$.}
{Fails for $Y=\varnothing$: the empty function from a one-point space of
measure zero into $\varnothing$ has no total extension.
\emph{Correction:} assume $Y\neq\varnothing$.\newline
\rocq{empty_target_extension_counterexample},
\rocq{measurable_function_total_extension}}
\entry{2.3.6, p.~77}{Corollary of Lusin's theorem: in case $X$ is
countably $\phi$ measurable, $f$ is $\phi$ almost equal to some Borel function
mapping $X$ into~$Y$.}
{Fails for $Y=\varnothing$ and a nonempty $X$ of measure zero.
\emph{Correction:} assume $Y\neq\varnothing$ or $X=\varnothing$.\newline
\rocq{borel_version_empty_target_counterexample},
\rocq{lusin_borel_version_metric}, \rocq{lusin_borel_version_radon}}
\entry{2.7.3, p.~123}{The sets $\bigcup_gA_g(V)\times A_g(V)$, for nonempty
open $V$, form a basis of an invariant uniformity.}
{Fails for $X=\varnothing$: the family is empty, whereas a uniformity
contains $X\times X$. \emph{Correction:} assume $X\neq\varnothing$. The
equivalence in the first assertion holds also for $X=\varnothing$.\newline
\rocq{empty_space_orbit_family_not_basis_example},
\rocq{uniform_action_orbit_basis}}
\entry{2.7.4, p.~123}{There exists a nonzero monotone invariant integral on
$\mathscr K(X)$.}
{Fails for $X=\varnothing$: then $\mathscr K(X)=\{0\}$, although all other
hypotheses hold. \emph{Correction:} assume $X\neq\varnothing$.\newline
\rocq{invariant_integral_empty_space_example},
\rocq{uniform_action_invariant_integral_exists}}
\entry{2.10.28(4), p.~192}{If $J\in F$, $2\le n\in\mathbf Z$,
$h(s)=\zeta(J)/n$ and $t=(\diam J-s)/(n-1)$, then $0<s<\diam J/n$, $t>s$,
$\Phi(J,n)$ is special for $J$, and an estimate holds.}
{Fails when $J$ is a point: then $\zeta(J)=h(0)=0$, hence $s=0$ and
$t=0$.
\emph{Correction:} assume $\diam J>0$.\newline
\rocq{zero_parent_cantor_split_impossible},
\rocq{is_equal_cost_interval_family}}
\entry{2.10.38, pp.~198--199}{Every compact $S\subset\R^n$ has points
$a_i,b_i\in S$ and an enclosing box with decreasing side lengths $c_i$ such
that $|(a_1-b_1)\wedge\cdots\wedge(a_m-b_m)|=c_1\cdots c_m$ for
$m=1,\ldots,n$.}
{Fails for $S=\varnothing$ and $n>0$: no points $a_1,b_1$ exist.
\emph{Correction:} assume $S\neq\varnothing$.\newline
\rocq{empty_compact_enclosing_box_endpoint_counterexample},
\rocq{federer_2_10_38}}
\entry{2.10.44, p.~202}{For $S=\{(1,-1),(-1,1),(1,1)\}$, the supremum norm
$\mu$ and the Euclidean norm $\nu$, $\mu(u-v)=2=\nu[f(u)-f(v)]$ for
$u,v\in S$.}
{Each point has distance $0$ from itself. \emph{Correction:} distinct
points have distance~$2$.\newline
\rocq{kirszbraun_printed_two_needs_distinct},
\rocq{kirszbraun_corner_pair_distances}}
\entry{2.10.44, p.~202}{If $S$ is a subset of a metric space $X$, every
Lipschitzian $f:S\to\R$ has a Lipschitzian extension $g:X\to\R$ with
$\operatorname{Lip}(g)=\operatorname{Lip}(f)$; one may define
$g(\xi)=\inf\{f(x)+\operatorname{Lip}(f)\operatorname{dist}(x,\xi):x\in S\}$.}
{For $S=\varnothing$ and $X\neq\varnothing$ the infimum is $\infty$.
\emph{Correction:} assume $S\neq\varnothing$ for the formula; the
extension theorem holds for $S=\varnothing$ with $g=0$.\newline
\rocq{empty_source_no_finite_mcshane_envelope},
\rocq{federer_2_10_44_scalar_formula}, \rocq{federer_2_10_44_scalar}}
\end{errors}

\subsection{Chapter 3}

\begin{errors}
\entry{3.1.1, p.~209}{$\limsup_{x\to a}|f(x)-f(a)|/|x-a|
=\sup\{|\langle v,D f(a)\rangle|:|v|=1\}=\|D f(a)\|$, and the
corresponding formula for $\liminf$ and $\inf$.}
{Fails for $X=\{0\}$: there is no $x\neq a$ and no $v$ with $|v|=1$, so
by the conventions of 2.1.1 (p.~51) the first two expressions are
$\sup\varnothing=-\infty$, whereas $\|D f(a)\|=0$.
\emph{Correction:} assume $X\neq\{0\}$. The definition of the differential and the closed unit ball
formula of the norm include this case.\newline
\rocq{norm_ratio_zero_source_exception},
\rocq{norm_differential_sharp_limits}}
\entry{3.2.16, p.~253}{If $f$ is $(\phi,m)$ approximately differentiable
at $a$, with $g$ as in the definition, then $f$ is differentiable relative
to $\{x:f(x)=g(x)\}$ at~$a$.}
{The concept of 3.1.22 (p.~235) is introduced under the standing assumption
$a\in\Clos S$, which 3.2.16 does not ensure. For $\phi=0$ and the empty
function every point is a point of approximate differentiability, while
every agreement set is empty. Without the standing assumption the
definition of 3.1.22 is satisfied vacuously, so the defect is a missing
presupposition, not a false conclusion. \emph{Correction:} add
$a\in\Clos\{x:f(x)=g(x)\}$, which holds when $\phi\{a\}>0$.\newline
\rocq{measure_approximate_empty_zero},
\rocq{measure_approximate_empty_agreement_not_closure},
\rocq{differentiable_representative_relative_agreement}}
\end{errors}

\subsection{Chapter 4}

\begin{errors}
\entry{4.1.1, p.~344}{The inclusion of $\mathscr D(U,Y)$ into
$\mathscr E(U,Y)$ is not a homeomorphic embedding unless $U=\varnothing$
or $Y=\{0\}$.}
{Fails when $X=\R^0$, $U=X$ and $Y=\R^1$. The source and target are
Banach spaces, $U$ is nonempty and $Y$ is nonzero. The whole domain is
compact, so $\mathscr D(U,Y)=\mathscr E(U,Y)$ and the two topologies agree.
\emph{Correction:} add $X\neq\{0\}$ to the printed assertion. More
generally, a noncompact open $U$ and nonzero $Y$ suffice. The actual
zero-dimensional counterexample and the corrected theorem are formalised.
This correction is absent from the published erratum.\newline
\rocq{zero_dimensional_test_topology_counterexample},
\rocq{smooth_test_nonzero_source_nonembedding},
\rocq{smooth_test_noncompact_nonembedding}}
\entry{4.1.25, p.~383}{For a compactly supported $\Leb^m$ summable
$m$-vectorfield $\xi$ and a Lipschitzian $f:\R^m\to\R^n$, the density
$\eta(y)=\sum_{x\in f^{-1}\{y\}}[\bigwedge_m\mathrm Df(x)]\xi(x)/J_mf(x)$
is a simple $m$-vector for $\Haus^m$ almost all $y$.}
{Fails for $m=n=0$. The empty product $1$ is the only simple
$0$-vector, so $0$ is not simple. The zero field on $\R^0$ and the
constant map give $\eta(0)=0$, and $\Haus^0$ is counting measure, which
has no nonempty null set. \emph{Correction:} assume $m>0$ for the
simplicity clause. The associated-tangent alternative and the signed
full-fibre formula hold in every degree.\newline
\rocq{current_area_zero_degree_simplicity_counterexample},
\rocq{current_lipschitz_area_geometry}}
\entry{4.1.27, p.~384}{For $u_0,\dots,u_m\in\R^n$, the map
$f(x)=u_0+\sum_{i=1}^m(\prod_{j=1}^ix_j)(u_i-u_{i-1})$ is univalent on
$C=\R^m\cap\{x:0<x_i<1\text{ for }i=1,\dots,m\}$ provided no $m-1$
dimensional affine subspace of $\R^n$ contains $\{u_0,\dots,u_m\}$.}
{Fails for $n=0$ and $m=2$. The space $\R^0$ has no $1$ dimensional
affine subspace, so the hypothesis holds vacuously. The map $f$ is
constant, and the distinct points $(\tfrac12,\tfrac12)$ and
$(\tfrac14,\tfrac14)$ of $C$ have the same image. \emph{Correction:}
assume $m\le n+1$, or assume that $u_0,\dots,u_m$ are affinely
independent. For $0<m\le n+1$ the printed hypothesis is equivalent to
linear independence of the edges $u_i-u_0$, and univalence holds.\newline
\rocq{simplex_cube_printed_univalence_counterexample},
\rocq{simplex_avoids_planes_nondegenerate},
\rocq{simplex_cube_printed_univalence}}
\end{errors}

\section{Convention-dependent qualifications}\label{s:conventions}
\noindent The first entries concern the zero exponent in dimension one.
The formalisation uses $0^0=1$; the following are not claims that the book
explicitly assigns that value in these formulas. All cases $n\ge2$ are
proved as printed, including zero volumes and empty cores. The later entries concern empty analytic germs, proper prime ideals,
and dimensions defined only for nonempty germs.
\subsection{Chapter 3}

\begin{errors}
\entry{3.2.42, p.~278}{The outer and inner distance-strip volumes,
divided by $r>0$, are bounded below by
\[
 n\alpha(n)^{1/n}v^{(n-1)/n},
\]
where $v$ is respectively the closure
volume or the volume of the inner core
$T_r=\{x:\operatorname{dist}(x,\R^n\sim S)\ge r\}$.}
{For $n=1$, the exponent is zero. With $0^0=1$ the inner formula fails
for $S=(0,1)$ and $r=2$: $T_2=\varnothing$, and the inner strip volume
divided by two is at most $1/2$, whereas the right side is
$\alpha(1)=2$. \emph{Qualification:} in dimension one retain a nonempty
source for the outer bound and a nonempty core for the inner bound.
The formalisation proves both such bounds, and both unrestricted bounds
for $n\ge2$.\newline
\rocq{boundary_layer_empty_core_line_exclusion_example},
\rocq{outer_boundary_layer_volume_nonempty},
\rocq{inner_boundary_layer_volume_nonempty_core},
\rocq{federer_3_2_42}}
\entry{3.2.43--44, p.~278}{The lower $(n-1)$-dimensional Minkowski
content of $\operatorname{Bdry}S$ is at least
$n\alpha(n)^{1/n}v^{(n-1)/n}$, where
$v=\Leb^n(\Clos S)<\infty$. For an $(n-1)$ rectifiable boundary,
3.2.44 replaces its lower content by its Hausdorff measure.}
{For $n=1$ and $S=\{0\}$, $v=0$ and the boundary lower degree-zero
content is one. Substituting $0^0=1$ would demand at least two.
\emph{Qualification:} retain positive closure volume in the line case;
all finite closure volumes, including zero, remain in every $n\ge2$.
The same qualification is retained in the rectifiable-boundary
corollary. No convention assigning a value to $0^0$ is inferred from
these printed inequalities.\newline
\rocq{isoperimetric_null_line_singleton_exclusion_example},
\rocq{positive_closure_volume_isoperimetric},
\rocq{federer_3_2_43}, \rocq{federer_3_2_44}}
\entry{3.4.5, pp.~320--321}{A variety germ is irreducible when it has no union decomposition into two proper variety subgerms. The book identifies this with primeness of its ideal and asserts a unique finite antichain of irreducible components for each germ under the Noetherian hypothesis.}
{The literal union test holds for the empty germ. Its vanishing ideal contains one, so it is not prime under the convention that prime ideals are proper. Moreover the empty family and the singleton consisting of the empty germ are distinct antichains of literal irreducibles with the same empty union. \emph{Qualification:} require irreducible component germs to be nonempty. Proper primeness is then equivalent to irreducibility, and the Noetherian decomposition is unique, with empty component family for the empty germ. The literal union test and its product-ideal characterization are also proved separately.\newline
\rocq{analytic_empty_germ_not_proper_prime},
\rocq{analytic_empty_germ_two_literal_families},
\rocq{analytic_irreducible_proper_prime_equivalence},
\rocq{noetherian_irreducible_component_decomposition}}
\entry{3.4.8(11),(16), pp.~327--328}{The boundary and analytic-hull formulas use the dimension of a variety germ.}
{Definition (3), p.~326, assigns dimension only to nonempty germs.
The zero-dimensional stratification has an empty boundary, and a
relatively open set in (16) may be empty. \emph{Qualification:} state
an empty germ as a separate alternative; apply the dimension bound
only in the nonempty case. No dimension of the empty germ is inferred
from the printed definition, and no contradiction is claimed for a
convention assigning it dimension $-\infty$. The formalisation retains
both alternatives explicitly.\newline
\rocq{analytic_point_stratification},
\rocq{federer_3_4_8_16}}
\entry{3.4.12: adaptation of 3.4.8(10)}
{p.~339, paragraph on replacing the ordered-sheet proof of (10);
p.~327, componentwise projection assertion}
{Replace the componentwise projection-homeomorphism conclusion by a
finite covering with local holomorphic sections. A connected complex
covering can have more than one point in each fibre. The polynomial graph
\[
 Z=\{(w,z)\in\mathbf C^2:w=z^2,\ z\ne0\}
\]
is connected and has exactly two points over every $w\ne0$. Its projection
cannot be injective on its one component. The actual local sheets,
their disjointness, exhaustive fibres and relatively open graph images
are proved in \texttt{ComplexSquareGraphDoubleCover.v}; the original
prime-model projection theorem uses the proved finite holomorphic covering
interface. The graph example checks the covering and component obstruction;
it does not separately instantiate its prime ideal and $E_m$ parameters.
This is an additional project counterexample and qualification,
not a published erratum entry.}
\end{errors}

\section{Misprints in statements, formulas, and the index}\label{s:misprints}
\subsection{Chapter 1}

\begin{errors}
\entry{1.9.1, p.~41}{$\bigodot_*V=\bigoplus_{m=2}^\infty\bigodot_mV$.}
{Read $m\ge0$, as the next sentence (degrees zero and one) requires. The
sum for the ideal starts correctly at $2$.\newline
\rocq{symmetric_zero_ambient_empty_word}}
\entry{1.9.4, p.~43}{$\Upsilon(e^\alpha)=\sum_{p=0}^m\sum
\binom{\alpha_1}{\beta_1}\cdots\binom{\alpha_m}{\beta_m}e^\beta\otimes
e^{\alpha-\beta}$ for $\alpha\in\Xi(n,m)$, the inner sum over
$\alpha\ge\beta\in\Xi(n,p)$.}
{The multi-indices have $n$ entries; read
$\binom{\alpha_1}{\beta_1}\cdots\binom{\alpha_n}{\beta_n}$.\newline
\rocq{symmetric_diagonal_monomial}}
\entry{1.10.6, p.~49}{$\alpha(j)=\card\{i:s(i)=j\}$ for $j=1,\ldots,m$.}
{Read $j=1,\ldots,n$.\newline
\rocq{symmetric_prescribed_fibre_condition},
\rocq{symmetric_fibre_cardinality}}
\end{errors}

\subsection{Chapter 2}

\begin{errors}
\entry{2.7.16(7), p.~141}{$s[p^{-1}\{y\}]=(p\circ s^{-1})^{-1}\{y\}$ for
$s\in\mathbf O(n)$, $y\in\R^n$.}
{Read $y\in\R^{n-m}$, since $p$ maps $\R^n$ into $\R^{n-m}$.\newline
\rocq{affine_parametrisation_fibre}}
\end{errors}

\subsection{Chapter 3}

\begin{errors}
\entry{3.2.22, p.~260}{$\int_W\operatorname{ap}J_\mu f\,d\Haus^m=\int_Z
\Haus^{m-\mu}(W\cap f^{-1}\{z\})\,d\Haus^mz$.}
{Read $d\Haus^\mu z$, as in assertion (3) on p.~258 and in the preceding
integrals.\newline
\rocq{rectifiable_coarea_subset_formula}, \rocq{federer_3_2_22_3}}
\entry{3.4.12: dimension comparison}{p.~339, final dimension display}
{The factor two is reversed. Read
\[
  \dim_{\mathbf R}\alpha=2\dim_{\mathbf C}\alpha,
\]
where the two dimensions are defined independently by the original
real and complex vanishing-ideal depths. The statement retains dimension
zero and every finite ambient dimension. It is proved in
\texttt{HolomorphicRealGermDimension.v}. This correction is not listed
in the May 1987 erratum.}
\entry{Index, p.~673}{The entry for Gross measure points to 2.10.3,
p.~172.}
{Read 2.10.4(1), p.~172.\newline
\rocq{gross_measure}}
\end{errors}

\subsection{Chapter 4}

\begin{errors}
\entry{4.1.8, p.~360}{The line after the definition of $\phi_S$
places $\eta$ in $\Lambda_j\R^m$.}
{Read $\eta\in\Lambda_j\R^n$. The vector belongs to the second
factor, and the preceding injection $Q$ has source $\R^n$.
The formal slice retains both original ambient dimensions and every
original hypothesis. This is a type correction, absent from the
repository transcription of the published erratum.\newline
\rocq{is_current_cartesian_slice_test},
\rocq{current_cartesian_slice_test_pairing}}
\entry{4.1.15, p.~373, final sentence}{The currents
$(\psi_i)_\#T$, where $T=\delta_0\wedge e_1\in\mathscr D_1(\R)$,
do not converge in $\mathscr D_0(\R)$.}
{Read $\mathscr D_1(\R)$. Smooth pushforward preserves current
degree, so these images are one-currents. The formalisation constructs
actual compact smooth scalar tests $\psi_i$ tending uniformly to zero
with one common differential bound and $\psi_i'(0)=(-1)^i$.
Their images are $(-1)^iT$ and have no weak limit in
$\mathscr D_1(\R)$. The pushed boundaries separately give a true
nonconvergence statement in $\mathscr D_0(\R)$. This correction is
absent from the published erratum.\newline
\rocq{point_one_current_oscillating_images},
\rocq{point_one_current_oscillating_images_diverge},
\rocq{point_one_current_oscillating_boundary_images_diverge}}
\entry{4.1.16, p.~373, line~29}{The published ambient-space
replacement in the nearest-point assertion.}
{The published erratum, p.~6, book p.~373, line~29, specifies
$\R^n$. The supplied book page already displays this symbol.
The formalisation retains arbitrary nonempty
closed convex subsets of this ambient space, including unbounded
sets and zero ambient dimension. This published correction is
separate from the local degree correction in the preceding entry.\newline
\rocq{is_euclidean_convex_nearest_map},
\rocq{euclidean_convex_nearest_map_graph},
\rocq{federer_4_1_16}}
\end{errors}

\section{Proof errors}\label{s:proofs}
\noindent The statements of these items are correct; an intermediate
assertion of the printed proof is false.
\subsection{Chapter 2}

\begin{errors}
\entry{2.10.24, p.~188}{Remark after the proof: the closed balls in a
boundedly compact metric space $X$ form a closed set with respect to the
Hausdorff distance.}
{False. In $X=[-1,-1/2]\cup[1/2,1]\cup\{-2,2\}\subset\R$ the set
$C=[-1,-1/2]\cup[1/2,1]$ is not a ball, but the balls with centre $1/2+h$
and radius $3/2-2h$ lie within Hausdorff distance $3h$ of $C$.
\emph{Correction:} the spherical equality is proved by Hausdorff
comparison instead.\newline
\rocq{closed_ball_hausdorff_closedness_counterexample},
\rocq{spherical_weighted_proper_equality}}
\entry{2.10.31, p.~196}{After the reflected-cap formula, $T_k\subset\B(a,s)$
for some $s<r$.}
{Here $r=\inf\{s:T\subset\B(0,s)$ for some $T\in C\}$ and $a$ lies on the
sphere $\operatorname{Bdry}\B(0,r)$, so the contradiction requires a ball
with centre $0$. The origin singleton lies in the centred unit ball and
misses its sphere, but lies in no ball of radius less than $1$ centred at a
point $a$ with $|a|=1$. \emph{Correction:} read $T_k\subset\B(0,s)$ for some
$s<r$.\newline
\rocq{shrinking_ball_wrong_center_counterexample},
\rocq{steiner_minimal_radius_contains_sphere}}
\end{errors}

\subsection{Chapter 3}

\begin{errors}
\entry{3.1.24, p.~241}{The norms of all terms of the coordinate series
total at most $\sum_{m=0}^\infty c^mm!\,(\sum_{i=1}^n|x_i-a_i|)^m/m!$;
the term of order $0$ is thus bounded by~$1$.}
{Condition (1) bounds only the differentials of positive order; the constant
analytic function $2$ violates the printed bound at the centre.
\emph{Correction:} keep $|f(a)|$ as a separate term. Absolute convergence
is unaffected.\newline
\rocq{analytic_printed_constant_term_bound_fails},
\rocq{analytic_coordinate_zero_degree_cost}}
\entry{3.2.29, p.~268}{$(\im F)\sim\bigcup\{f(U):U\in G\}
=F(\R^m\sim C)\cup f(C\sim P)$.}
{Only the inclusion $\subset$ holds: $F(\R^m\sim C)$ may meet
$\bigcup\{f(U):U\in G\}$. \emph{Correction:} replace $=$ by $\subset$,
which suffices for the subsequent estimate.\newline
\rocq{agreement_image_cover_remainder},
\rocq{lipschitz_small_manifold_remainder}, \rocq{federer_3_2_29}}
\entry{3.2.38, p.~275}{The first displayed Steiner sum in the bound
for $\rho(K)$ runs from $j=1$ to $n$.}
{It omits the constant intrinsic-volume term. For the actual zero affine
map on the unit line cube, $m=n=1$, the image is $\{0\}$ and its
normalized radius-$r$ tube volume is $2r>0$, while the only printed
$j=1$ coefficient is zero. \emph{Correction:} sum from $j=0$ to $n$.
The extra term is absorbed in the unspecified finite error constant;
the statement of the lemma is unchanged.\newline
\rocq{missing_steiner_constant_counterexample_exists},
\rocq{compact_minkowski_constant_map_correction_example},
\rocq{affine_coordinate_cube_steiner_estimate}}
\entry{3.3.19, p.~305}{The displayed difference of closed caps lies in
$X(a,2rs,V,s)$ when $\sin\alpha=s$, with the original strict radius.}
{This auxiliary inclusion fails at every actual parent level $j\ge1$ on
the fixed boundary branch. Take $a=\exp(0i)$, $r=\rho_j$,
$\alpha=\pi/6$, $s=1/2$ and $x=z+r\exp(i\pi/3)$.
The first closed cap contains $x$, while its radial coordinate $r/2$ is
strictly below $r(1-2/(j+4))$, so $x$ is outside the second cap.
But $|x-a|=r=2rs$, which the original strict cone excludes.
\emph{Repair:} choose an auxiliary sine aperture $0<\sigma<s$ and a
local normal buffer, then take the original parent index sufficiently
large to absorb the radial deficit $2r/(j+4)$. Letting these finite
buffers decrease proves the printed $1/(2\pi)$ lower conical density
with no new numbered-statement hypothesis. The main assertion is
unchanged; this auxiliary correction is absent from the published erratum.\newline
\rocq{circular_disc_actual_cap_counterexample},
\rocq{circular_disc_buffered_cone_mass},
\rocq{circular_disc_conical_upper_density_almost}}
\entry{3.3.20, p.~308}{Every $S\subset\R^2$ has
$\diam S\ge\psi_j(S)$ for some finite original stage $j$.}
{Qualify this auxiliary finite-stage assertion by $0<\diam S<\infty$.
The original corner singleton $S=\{(0,0)\}$ has diameter zero
but strictly positive $\psi_j(S)$ at every finite original stage, since
an original level square always contains that corner. Its actual product
measure is zero. The positive-diameter finite-stage argument, together
with this separate diameter-zero product nullity, gives every original
Hausdorff-generator bound and preserves the numbered conclusion
$\Haus^1(E)=\sqrt2$. This auxiliary qualification is absent from the
published working erratum.\newline
\rocq{critical_square_original_singleton_qualification},
\rocq{critical_original_positive_diameter_attainment},
\rocq{critical_original_cantor_square_sharp_length}}
\end{errors}

\section{An error and a normalization issue in the list of corrections
of May 1987}\label{s:erratum}
\subsection{Chapter 2}

\begin{errors}
\entry{2.9.18(1), p.~163}{Erratum p.~3, book p.~163 line~11: in
$\psi(A)=\sup(\{0\}\cup\{c_j:A\cap S_j\neq\varnothing\})$ read
$A\cap S_j=\varnothing$.}
{With this replacement $\psi(\varnothing)=c_1=1$, so $\psi$ is not a
measure, and the computation $\psi\B(x,r)=c_{j+1}$ for
$c_{j+1}\le r\le jc_{j+1}$ fails. \emph{Correction:} keep the book's
condition $A\cap S_j\neq\varnothing$.\newline
\rocq{weighted_miss_capacity_empty_positive},
\rocq{weighted_miss_capacity_not_measure},
\rocq{weighted_hit_measure_is_measure}}
\end{errors}

\noindent\textbf{Normalization note for 3.2.38 (editorial).}
On erratum p.~5, book p.~275 line~2, the replacement is $r^{n-j}$.
The original scan confirms this exponent. In the displayed bound for
$\rho(K)$, however, the tube volume has already been divided by
$\alpha(n-m)r^{n-m}$ (definition on book p.~274). A raw Steiner term has
power $r^{n-j}$; its normalized power is therefore $r^{m-j}$, as printed
in the book. The formalisation keeps the full raw Steiner polynomial
and divides by the actual denominator, retaining $j=0$ as above; see
\rocq{affine_coordinate_cube_steiner_estimate}.
Thus the erratum's replacement cannot be applied literally to that
already normalized power. This is an algebraic consistency observation;
no separate Rocq counterexample to this erratum replacement is claimed.
The transcription \texttt{federer-erratum.tex} preserves the scan.

\section{Corrections of the list of May 1987 confirmed by the
formalisation}\label{s:confirmed}
\noindent The formalisation adopts the following corrections of Federer's
list. Pages and lines are those of the book. The table includes corrections
used in the formalised statements or arguments; purely stylistic and
bibliographical changes are omitted.
\subsection{Chapter 1}

\begin{confirmed}
\entry{1.5.1}{22, 9}{In the associativity identity of interior
multiplication the multivector is $\zeta\in\bigwedge_tV$.}
\entry{1.7.8}{34, 29--30}{$\langle\xi,*\beta\rangle=(-1)^{p(n-p)}\langle
*\xi,\beta\rangle$ for $\xi\in\bigwedge_pV$, $\beta\in\bigwedge^{n-p}V$.}
\entry{1.7.9}{35, 3}{$(v\otimes w)\bullet(v'\otimes w')=(v\bullet
v')(w\bullet w')$.}
\end{confirmed}

\subsection{Chapter 2}

\begin{confirmed}
\entry{2.1.5}{57, 26}{$\phi(A)=\card(A)$ for $A\subset S$.}
\entry{2.2.1}{60, 4}{Both generation criteria assume $\varnothing\in S$;
for $S=\varnothing$ the generated family $\{\varnothing\}$ is not a Borel
family over a nonempty $X$.}
\entry{2.2.10}{66, 13 and 15}{A sequence $S_1,S_2,S_3,\ldots$ of closed
sets in the cylinders $\mathscr N_j$.}
\entry{2.2.10}{67, 2}{$\bigcup G$ and $\bigcup H$ are Borel separated.}
\entry{2.3.8}{78, 7}{$\mathbf A(\phi,Y)$ is a set of partial functions; the
vector space and convex hull statements concern its members with domain
$X$.}
\entry{2.4.12}{86, 26--27 and bottom; 87, 1}{$\phi_{(p)}$ is defined on
$\mathbf A(\phi,Y)$ and is a seminorm on $\widetilde{\mathbf A}(\phi,Y)$.}
\entry{2.4.12}{87, 18}{The functions with finite image are dense in
$\Lsp_p(\phi,Y)$ provided $p<\infty$ or $\dim Y<\infty$.}
\entry{2.4.12}{87, 32; 88, 3}{$\phi_{(1)}$ closed subspace; the subspace of
functions with finite image.}
\entry{2.4.16}{89, 19}{The tests range over $u\in\Lsp_q(\phi)$ with
$\phi_{(q)}(u)\le1$.}
\entry{2.5.1}{91, 28}{$\sup\{f,g\}\in L$ for all $f,g\in L$.}
\entry{2.5.7}{98, 10--11}{$\mu$ is defined on $\Lt_p(\phi)$, and the
supremum is over $\phi_{(p)}(f)\le1$.}
\entry{2.5.7}{99, 13, 14, 17}{The residual Daniell integral is taken on
$L^+$.}
\entry{2.5.8}{100, 3 and 5}{The Daniell integral is on $\Lt_\infty(\phi)$
and vanishes whenever $\phi_{(\infty)}(f)=0$; without this hypothesis
evaluation at a point under the zero measure is a counterexample.}
\entry{2.5.9}{100, 21--24}{$\Lt_1(\phi)$ in place of $\Lsp_1(\phi)$;
$L^+$ is dense in $\Lt_1(\phi)^+$.}
\entry{2.5.11}{102, 8 and 12}{The lattice is
$\Lt_1(\alpha_1\times\beta_2)\cap\Lt_1(\beta_1\times\alpha_2)$, and
$\Lt_1(\phi)=L$.}
\entry{2.6.4}{118, 20, 21, 24; 119, 3, 4, 6}{$\Lt_1$ in place of
$\Lsp_1$.}
\entry{2.6.4}{119, 19}{$\nu(h)=\int h\,d(\alpha\times\beta)$ for
$h\in\mathscr K(X\times Y)$.}
\entry{2.7.4}{124, 30; 125, 7--8}{``Whenever $0\neq\nu\in\mathscr K(X)^+$,
$(\spt\nu)\times(\spt\nu)\subset\beta$, $a\in\spt\nu$'', and
$\lambda\in S(\beta)$.}
\entry{2.7.14}{131, 10}{The measure associated with a \emph{nonzero}
monotone left (right) $\chi$ covariant integral is a Radon measure. For the
zero integral on the group of subsets of $\N$ under symmetric difference the
associated measure is not a Radon measure
(\rocq{zero_left_invariant_measure_not_radon}).}
\entry{2.7.14}{132, 22 and 25}{$f\in\Psi$ in place of $f\in\Phi$.}
\entry{2.7.14}{133, 3}{$U=(G\sim P)\cup(P\cap U)$.}
\entry{2.7.15}{133, 4}{The measure associated with a \emph{nonzero}
monotone $\chi$ covariant integral is a Radon measure
(\rocq{homogeneous_covariant_measure_radon}).}
\entry{2.7.16(1)}{134, 26; 135, 35}{$\Leb^n$ is the unique \ldots; the
exponent in the recursion for $\bm\alpha(n)$ is $(n-1)/2$.}
\entry{2.7.16(2)}{136, 29}{The left Haar integral of $GL(V)$ is $\int
u\,|\det|^{-n}\,d\phi$.}
\entry{2.7.16(5)}{138, 34}{The modular function of $G^W$ involves
$|\det[L(g)|W]|^{-1}$.}
\entry{2.7.18}{141, 25}{$\chi(g)\to\chi(h)$.}
\entry{2.8.2}{142, 33}{$G_0=\varnothing$.}
\entry{2.8.4}{143, 17}{$\varnothing\notin F$.}
\entry{2.8.17}{151, 24; 152, 10}{$\phi$ is positive on $V(X)$; cite
2.8.3.}
\entry{2.9.7}{155, 26}{$\mathbf D(\psi,\phi,V,x)\notin\R$.}
\entry{2.9.12}{159, 7}{$X\sim\{z:f(z)\le t\}$ has density $0$ at $x$.}
\entry{2.9.13}{159, 19}{$(\phi,V)$.}
\entry{2.9.15}{161, 11}{Cite 2.8.14.}
\entry{2.9.17}{162, 24--25}{$\B(x,2r)$.}
\entry{2.10.1}{170, 24 and 35}{Cite [F12].}
\entry{2.10.14}{177, 23}{The upper integral $\int^*$.}
\entry{2.10.21}{184, 8 and 10}{Nonempty compact convex subsets.}
\entry{2.10.28(3)}{192, 4}{$\card(H_0)=1$.}
\entry{2.10.28(4)}{193, 7}{Sum over $S\in\Phi(j,n)$ with $S\subset T$.}
\entry{2.10.29}{194, 8}{$\sigma_j\ge t_{j+1}\to0$.}
\entry{2.10.32}{197, 4}{$\mu\bm\alpha(m)^{-1}2^m$.}
\end{confirmed}

\subsection{Chapter 3}

\begin{confirmed}
\entry{3.1.1}{210, 26}{$u_{n+1}-u_n=g(u_{n-1})-g(u_n)$.}
\entry{3.1.4}{214, 24}{$\ge1/j$.}
\entry{3.1.9}{218, 12}{$\U(z,1/j)$.}
\entry{3.1.11}{219, 4}{$\{|\phi(v_1,\ldots,v_m)|:\ldots\}$.}
\entry{3.1.11}{221, 1}{$a+t(x-a)$.}
\entry{3.1.12}{224, 4}{$\alpha h(x)+(\beta+1)h(s)$.}
\entry{3.1.16}{228, 19}{All $a$ in $A$.}
\entry{3.1.20}{232, 11, 14--15}{Maps $f_i:W_i\to B$ of class $k$ with
$f_i(x)=x$ for $x\in B\cap W_i$; the blend maps into $T\cap B$.}
\entry{3.1.22}{235, 19}{$f(a)=\eta$ in case $a\in S$.}
\entry{3.2.2}{242, 23}{$\mathbf O(n,m)$.}
\entry{3.2.3}{244, 11}{$J_mf(x)$.}
\entry{3.2.6}{246, 3 and 13}{$f(a)\le f(b)$; $\Leb^1$ summable.}
\entry{3.2.10}{248, 21 and 25}{$\Leb^{m-n}$ and $\Haus^{m-n}$.}
\entry{3.2.11}{249, 21}{$\Leb^m(A)$.}
\entry{3.2.16}{253, 29}{$\delta_{i+1}<\delta_i/2$.}
\entry{3.2.16}{253, 32}{Differentiable relative to $T$ at $a$; $T\cup\{a\}$
in case $\phi\{a\}>0$.}
\entry{3.2.16}{254, 2}{``generalizes'' in place of ``coincides with''.}
\entry{3.2.21}{257, 26}{$\Leb^\mu(B)$.}
\entry{3.2.22}{258, 20 and 22; 259, 3--4}{The source remainder
$R=W\sim\bigcup_{i=1}^\infty\psi_i(K_i)$, with $\Haus^m(R)=0$ and
$\Haus^{m-\mu}$ null fibres.}
\entry{3.2.25}{261, 4}{$\Haus^m$.}
\entry{3.2.26}{261, 22 and 29; 262, 14}{Cite 2.10.5; $x-w$; cite 3.2.20.}
\entry{3.2.28(2)}{265, 3}{Delete ``because $\mathbf O(n)$ is compact,''.}
\entry{3.2.28(3)}{265, 12}{$\mathbf O(m,m-1)$.}
\entry{3.2.33}{270, 25 and 28}{Cite 3.1.24 for the analytic rank
alternative, and 3.1.18 for positivity of the complementary-dimensional
fibres.}
\entry{3.2.35}{271, 24}{Parametrize the convex boundary by
$x\mapsto g(x)^{-1}x$ on the unit sphere, where $g$ is the gauge.}
\entry{3.2.36}{273, 18}{Read ``$m$-th'' in the elementary symmetric
function description of the box coefficient.}
\entry{3.2.38}{275, 14}{Use the actual radius-$r$ tube
$\{z:\operatorname{dist}[z,f(C)]<r\}$ in the final volume quotient.}
\entry{3.3.2}{288, 1--6}{Use the strict margin $\lambda+1/j$ and
the union over $j$ of intersections over $i$ for the strict superlevel,
as in the published erratum pp.~5--6
(\rocq{structure_corrected_strict_superlevel}).}
\entry{3.3.20}{307, 7}{Delete the false assertion that
$c_u<1$ always implies $\Leb^1[p_u(E)]=0$.
Published erratum p.~6. The formalisation does not state or use it.}
\entry{3.3.20}{307, 26}{Use projected first-level lengths
$2(u_1/4+u_2)$ on $0\le u_2\le u_1/2$ and $3u_1/2$ on
$u_1/2\le u_2\le u_1$, and the universal bound $3/\sqrt5$.
Published erratum p.~6; both exact branches and every signed original
direction are proved (\rocq{critical_original_projected_union_length},
\rocq{critical_original_orthogonal_projection_bound}).}
\end{confirmed}

\subsection{Chapter 4}
\begin{confirmed}
\entry{4.1.10}{364, 11}{Replace $\mathscr D_j$ by $\mathscr D^j$, as on
erratum p.~6. The assertion concerns compact test forms, not currents.
The proof retains this test class and identifies the original smooth
homotopy primitive; it does not claim compact support for that primitive.
The same argument and all three original compact-current conclusions are
proved on every star-shaped open domain. See
\rocq{current_cone_test_poincare},
\rocq{current_cone_test_primitive_integral}, and
\path{docs/checkpoints/4.1.10-validation.json}.}
\entry{4.1.17}{374, 16}{Use $Q+\partial S$, as specified on
erratum p.~7. Both $Q+\partial S$ and $Q-\partial S$ are normal,
but the plus sign gives the required mass error
$\mathbf M(T-Q-\partial S)$. The formalisation proves the exact
normal mass-closure equality with the original compact index and
the contraction assertion for every variation-summable form.
No new local source error is claimed. See
\rocq{flat_finite_mass_normal_approximation},
\rocq{federer_4_1_17_mass_closure}, and
\rocq{federer_4_1_17_summable_contraction}.}
\entry{4.1.27}{384, 19}{The last factor of the wedge formula is
$\langle e_m,\mathrm Df(x)\rangle$, as on erratum p.~7. The
formalisation proves
$\langle e_1,\mathrm Df(x)\rangle\wedge\dots\wedge\langle e_m,\mathrm Df(x)\rangle
=(\prod_{i=1}^{m-1}x_i^{m-i})\zeta$ at every $x\in\R^m$, where
$\zeta=(u_1-u_0)\wedge\dots\wedge(u_m-u_{m-1})$. See
\rocq{simplex_cube_wedge_differential}.}
\end{confirmed}

\section{Further published corrections in Chapter 3}\label{s:further-chapter3}

These five entries reproduce corrections in Federer's May 1987 list, not new project discoveries.
The book passages were checked against the retained erratum transcription. The bibliographical insertion in 3.4.4
changes no mathematical statement. Item 3.2.45 is complete with the
project correction recorded above; its entries here retain the published
source corrections.
Both corrections to 3.4.8 are now applied in registered proofs, as
recorded below.

\subsection{Bibliographical references}
\begin{confirmed}
\entry{3.2.45}{279, 27}{Replace the mixed-volume reference marker by
\texttt{[EG 2,}. This is an editorial correction in erratum p.~5.
It is recorded without a claim to have formalised the cited external result.}
\entry{3.4.4}{318, 26}{Insert \texttt{[HF 2] or} before the existing
reference to \texttt{[W 2]}, as prescribed on erratum p.~6.
The supplementary arc examples remain bibliographical source accounts;
their separate constructions are not claimed formalised.}
\end{confirmed}

\subsection{Analytic stratification: newly validated published corrections}
\begin{confirmed}
\entry{3.4.8(1)}{326, 3}{The covector basis belongs to
$\bigwedge^1\R^n$, as on erratum p.~6. Retain the distinction between
vectors and covectors in the two dual orthonormal bases.}
\entry{3.4.8(11)}{327, 24}{Use the complete replacement on erratum p.~6:
\[
 S_i\cap\bigcup_{0\le j<i}S_j=\varnothing,
 \qquad V_i\cap\Clos S_i\subset S_i\cup G_i,
\]
where $G_i$ is an analytic subvariety of $V_i$ satisfying
\[
 \dim\gamma_0(G_i)<m,
 \qquad G_i\subset\bigcup_{0\le j<i}S_j.
\]
Both the dimension bound and containment in the earlier strata are part
of the correction. Every corrected law is proved in
\texttt{Federer348Stratification.v}. The compact-component obligation
of 3.4.5 is discharged in \texttt{AnalyticSetCompactComponents.v}.
Empty boundary germs are explicit because 3.4.8(3) defines dimension
only for nonempty germs. See \rocq{docs/checkpoints/3.4.8-validation.json}.}
\end{confirmed}

\subsection{Further mathematical corrections}
\begin{confirmed}
\entry{3.2.45}{280, 6}{Use
$\mathscr Q_\infty^m\le\mathscr Q_1^m$, as on erratum p.~5.
The consulted reprint already displays this inequality; this entry is
an erratum cross-reference, not a further error in that copy.
The sharp comparison is proved for every target set, including degree
zero and infinite masses, by
\rocq{convex_integral_geometric_infinity_below_first}.}
\end{confirmed}
\end{document}
